% Worked Examples: Concept 3.9.4 - Robust Pole Placement
\documentclass[11pt,a4paper]{article}
\usepackage{worked-examples-template}

\title{\textbf{Worked Examples}\\[0.5em]
\large Concept 3.9.4: Robust Pole Placement\\[0.3em]
\normalsize \coursesubtitle{} -- \coursetitle}
\author{Assoc.\ Prof.\ William Robertson \and Dr Sean McGowan}
\date{\today}

\begin{document}

\maketitle
\tableofcontents
\newpage

\section{Concept 3.9.4: Robust Control}

\subsection{Example 1: PI Controller Design with Fast Stable Pole}

When a process has stable poles faster than the desired closed-loop bandwidth, special care is needed to avoid robustness problems.

\paragraph{Problem:}
Consider a first-order process $P(s) = 2/(s+10)$ with a fast stable pole at $s = -10$. Design a PI controller $C(s) = k_p + k_i/s$ to place the closed-loop poles at $s = -2 \pm \ii$.

(a) Calculate the required controller gains $k_p$ and $k_i$.

(b) Determine the sensitivity function $S(s)$ and complementary sensitivity function $T(s)$.

(c) Find the maximum sensitivity $M_s = \max_\omega |S(\ii\omega)|$ and maximum complementary sensitivity $M_t = \max_\omega |T(\ii\omega)|$.

(d) Explain why this design has poor robustness and suggest an improvement.

\paragraph{Solution:}

\textbf{Step 1: Closed-loop characteristic equation}

The process transfer function is $P(s) = 2/(s+10)$ (where $b = 2$ and $a = 10$).

The PI controller is $C(s) = k_p + k_i/s = (k_ps + k_i)/s$.

The loop transfer function:
\[
L(s) = P(s)C(s) = \frac{2}{s+10} \cdot \frac{k_ps + k_i}{s} = \frac{2k_ps + 2k_i}{s(s+10)}
\]

The closed-loop characteristic equation is:
\[
1 + L(s) = 0 \quad \Rightarrow \quad s(s+10) + 2k_ps + 2k_i = 0
\]
\[
s^2 + (10 + 2k_p)s + 2k_i = 0
\]

\textbf{Step 2: Pole placement calculation}

For poles at $s = -2 \pm \ii$, the desired characteristic polynomial is:
\[
(s - (-2 + \ii))(s - (-2 - \ii)) = (s+2-\ii)(s+2+\ii) = (s+2)^2 + 1
\]
\[
= s^2 + 4s + 4 + 1 = s^2 + 4s + 5
\]

Matching coefficients with $s^2 + (10 + 2k_p)s + 2k_i = 0$:
\[
10 + 2k_p = 4 \quad \Rightarrow \quad k_p = -3
\]
\[
2k_i = 5 \quad \Rightarrow \quad k_i = 2.5
\]

\textbf{Step 3: Sensitivity functions}

The controller is $C(s) = -3 + 2.5/s = (-3s + 2.5)/s$.

Notice that $k_p < 0$, which means the controller has a zero in the right half-plane:
\[
C(s) = \frac{-3s + 2.5}{s} = -3\frac{s - 0.833}{s}
\]

The zero is at $s = 2.5/3 = 0.833$ (RHP zero).

Loop transfer function:
\[
L(s) = \frac{2(-3s + 2.5)}{s(s+10)} = \frac{-6s + 5}{s(s+10)}
\]

Sensitivity function:
\[
S(s) = \frac{1}{1 + L(s)} = \frac{s(s+10)}{s(s+10) + (-6s + 5)} = \frac{s(s+10)}{s^2 + 4s + 5}
\]

Complementary sensitivity:
\[
T(s) = \frac{L(s)}{1 + L(s)} = \frac{-6s + 5}{s^2 + 4s + 5}
\]

\textbf{Step 4: Calculate maximum sensitivities}

For $M_s$, we need to find $\max_\omega |S(\ii\omega)|$:
\[
S(\ii\omega) = \frac{\ii\omega(\ii\omega + 10)}{(\ii\omega)^2 + 4\ii\omega + 5} = \frac{-\omega^2 + 10\ii\omega}{-\omega^2 + 4\ii\omega + 5}
\]
\[
= \frac{(5-\omega^2) + 10\ii\omega}{(5-\omega^2) + 4\ii\omega}
\]

The magnitude is:
\[
|S(\ii\omega)| = \frac{\sqrt{(5-\omega^2)^2 + 100\omega^2}}{\sqrt{(5-\omega^2)^2 + 16\omega^2}}
\]

At the resonant frequency $\omega_r = \sqrt{5} \approx 2.24$ rad/s (where the denominator is minimized):
\[
|S(\ii\sqrt{5})| = \frac{\sqrt{0 + 500}}{\sqrt{0 + 80}} = \frac{\sqrt{500}}{\sqrt{80}} = \frac{22.36}{8.94} = 2.50
\]

Therefore, $M_s \approx 2.50$ (very large, indicating poor robustness!).

\textbf{Step 5: Analysis and improvement}

\textbf{Why is robustness poor?}
\begin{itemize}
\item The fast stable pole at $s = -10$ is much faster than the desired bandwidth ($\omega_n = \sqrt{5} \approx 2.24$)
\item To place slow poles at $s = -2 \pm \ii$, we need $k_p < 0$, creating an RHP zero in the controller
\item This RHP zero causes a large sensitivity peak: $M_s = 2.50 > 2$, indicating poor stability margins
\item The minimum robust stability margin is $1/M_s = 0.4$, meaning only 40\% multiplicative uncertainty is tolerable
\end{itemize}

\textbf{Improvement strategy:}
\begin{itemize}
\item Choose closed-loop poles faster than the process pole: e.g., $s = -15 \pm 5\ii$
\item This ensures $k_p > 0$ and avoids the RHP zero
\item Accept the natural fast dynamics of the process rather than fighting them
\item Alternatively, use a lead compensator to cancel the fast pole (if exact cancellation is feasible)
\end{itemize}

\textbf{Key insights:}
\begin{itemize}
\item Attempting to make the closed-loop slower than a fast stable process pole leads to poor robustness
\item Negative proportional gain creates RHP zeros that amplify sensitivity
\item Robust design: let fast stable poles remain fast in closed-loop
\end{itemize}

\subsection{Example 2: Pole Placement with Slow Stable Zero}

Process zeros slower than the desired bandwidth also require careful treatment to maintain robustness.

\paragraph{Problem:}
Consider the vehicle steering system with transfer function:
\[
P(s) = \gamma\frac{s + 1/\gamma}{s^2} = \frac{s + 0.2}{s^2}
\]
where $\gamma = 5$ m is a characteristic length and the zero is at $s = -0.2$ rad/s (slow).

A state feedback controller with observer is designed to place all closed-loop poles at $s = -5$. The resulting controller transfer function is:
\[
C(s) = \frac{125s + 250}{s^2 + 20s + 100}
\]

(a) Calculate the loop transfer function $L(s)$.

(b) Determine the complementary sensitivity function $T(s)$.

(c) Find the maximum of $|T(\ii\omega)|$ and identify if there is a robustness issue.

(d) Suggest a modification to improve robustness.

\paragraph{Solution:}

\textbf{Step 1: Loop transfer function}

\[
L(s) = P(s)C(s) = \frac{s + 0.2}{s^2} \cdot \frac{125s + 250}{s^2 + 20s + 100}
\]
\[
= \frac{(s + 0.2)(125s + 250)}{s^2(s^2 + 20s + 100)} = \frac{125(s + 0.2)(s + 2)}{s^2(s^2 + 20s + 100)}
\]

\textbf{Step 2: Complementary sensitivity function}

The complementary sensitivity is:
\[
T(s) = \frac{L(s)}{1 + L(s)} = \frac{P(s)C(s)}{1 + P(s)C(s)}
\]

The closed-loop poles are all at $s = -5$ (quadruple pole), so:
\[
1 + L(s) = \frac{(s+5)^4}{s^2(s^2 + 20s + 100)}
\]

Therefore:
\[
T(s) = \frac{125(s + 0.2)(s + 2)}{(s+5)^4}
\]

\textbf{Step 3: Analyze maximum of $|T(\ii\omega)|$}

The numerator has zeros at $s = -0.2$ and $s = -2$.
The denominator has a quadruple pole at $s = -5$.

At low frequencies ($\omega \ll 0.2$), the magnitude behaves as:
\[
|T(\ii\omega)| \approx \frac{125 \cdot 0.2 \cdot 2}{5^4} = \frac{50}{625} = 0.08
\]

At the slow zero frequency $\omega = 0.2$:
\[
T(\ii \cdot 0.2) = \frac{125(\ii \cdot 0.2 + 0.2)(\ii \cdot 0.2 + 2)}{(\ii \cdot 0.2 + 5)^4}
\]

This is complex to evaluate analytically, but we can observe that the slow zero at $s = -0.2$ will cause a peak in $|T|$ near $\omega = 0.2$ rad/s.

By numerical evaluation or frequency response analysis, the maximum occurs near $\omega \approx 0.2$ rad/s with $M_t \approx 3.5$ (very large!).

\textbf{Step 4: Robustness issue and improvement}

\textbf{Problem:}
\begin{itemize}
\item The slow stable zero at $s = -0.2$ is much slower than the closed-loop poles at $s = -5$
\item This creates a large peak in the complementary sensitivity: $M_t \approx 3.5$
\item From robust stability: $|\delta| < 1/M_t \approx 0.29$, tolerating only 29\% multiplicative uncertainty
\item The system has poor robustness to high-frequency unmodeled dynamics
\end{itemize}

\textbf{Improvement:}
Instead of placing all poles at $s = -5$, place one closed-loop pole near the slow zero:

Modified pole placement: poles at $s = -0.2, -5, -5, -5$

This creates a controller pole that cancels (or nearly cancels) the process zero:
\[
C_{\text{improved}}(s) = K\frac{(s + 5)^3}{s^2 + 20s + 100} \cdot \frac{1}{s + 0.2}
\]

The cancellation results in:
\[
T_{\text{improved}}(s) \approx \frac{K(s + 2)}{(s+5)^3}
\]

This eliminates the slow zero from $T$, reducing the complementary sensitivity peak to $M_t \approx 1.2$, greatly improving robustness.

\textbf{Key insights:}
\begin{itemize}
\item Slow stable zeros should be canceled by closed-loop poles to avoid peaks in $T$
\item Cancellation improves robustness to high-frequency uncertainty
\item Design rule: assign controller poles near slow process zeros
\end{itemize}
\end{document}
